\section{Division}

\begin{rappelbox}
    On appelle \bb{quotient} le résultat d'une \answer{division}. Pour a et b deux nombres
    relatifs, on note $a\div b$ ou $\dfrac{a}{b}$ le quotient de a par b.
\end{rappelbox}

\begin{propbox}
    Pour \bb{diviser} deux nombres relatifs: 
    \begin{itemize}
        \item On divise leurs valeurs absolues (ou distances à zéros);
        \item On applique \bb{la même règle des signes que pour la multiplication} :
            \begin{itemize}[label=\textbullet]
                \item le quotient de deux nombres de \bb{même signe} est \bb{positif} ;
                \item le quotient de deux nombres de \bb{signes différents} est \bb{négatif}
            \end{itemize}
    \end{itemize}
\end{propbox}

\begin{exemplebox}
    $$ (-12)\div (-3)=\answer{4} $$
    $$ 12\div (-3)=\answer{-4} $$
\end{exemplebox}

\begin{exercicebox}
    Calculer les quotients suivants:
    $$18\div (-6) = \answer{-3}$$
    $$(-49)\div (-7) = \answer{7}$$
    $$(-56)\div 7 = \answer{-8}$$
\end{exercicebox}

\begin{propbox}
    Soient a et b deux nombre relatifs avec b \neq 0. Alors
    \begin{enumerate}
        \item $\dfrac{-a}{-b} = \dfrac{a}{b}$
        \item $\dfrac{a}{-b} = \dfrac{-a}{b} = -\dfrac{a}{b}$
    \end{enumerate}
\end{propbox}
